\chapter{Introduction}
\label{chap:introduction}

% YOU WRITE THIS CHAPTER. Notes below are structure/content pointers only,
% not sentences to copy. Ask me to explain any of these before you write
% about them.

\section{Overview}

Large Language Models (LLMs) have become increasingly capable of solving
complex reasoning tasks, including mathematical problem solving, making
them valuable for applications such as intelligent tutoring systems,
educational technologies, and decision-support tools. As these models are
increasingly used in domains that require reliable multi-step reasoning,
understanding not only their final performance but also the reasoning
process leading to their answers has become an important area of research.

This dissertation investigates the mathematical reasoning performance of
three Large Language Models---Gemini Flash Lite, GPT-OSS-120B, and Mistral
Large---using two reasoning-based prompting strategies: Chain-of-Thought
(CoT) and Tree-of-Thought (ToT). The study evaluates model performance on a
curated dataset of mathematical problems drawn from the Hendrycks MATH
dataset and selected American Invitational Mathematics Examination (AIME)
questions. Beyond measuring answer correctness, the research analyses the
types of mathematical reasoning errors produced by each model, identifies
where within the reasoning process these errors occur, and compares the
effectiveness and computational cost of the two prompting strategies.

This chapter introduces the research by presenting the motivation for the
study, defining the research problem, outlining the research questions and
objectives, and summarising the main contributions of the dissertation. It
concludes with an overview of the remaining chapters, providing the
context for the methodology, experimental evaluation, statistical
analysis, and discussion presented throughout the dissertation.

\section{Motivation}

Large Language Models (LLMs) have become increasingly capable of solving
complex reasoning tasks and are now widely used in applications such as
intelligent tutoring systems, educational technologies, and
decision-support tools. As these models continue to be integrated into
domains that require reliable mathematical reasoning, understanding the
quality of their reasoning process has become increasingly important.
Mathematical problem solving differs from many other natural language
tasks because it requires logically consistent multi-step reasoning, where
an error in an intermediate step can invalidate the entire solution.
Consequently, evaluating the reasoning capabilities of LLMs is essential
for assessing their reliability and suitability for real-world
applications.

The growing use of reasoning-based prompting techniques has created new
opportunities to improve the mathematical performance of LLMs.
Chain-of-Thought (CoT) prompting encourages models to generate intermediate
reasoning steps, while Tree-of-Thought (ToT) prompting extends this process
by exploring multiple reasoning paths before selecting a final solution. As
these prompting strategies become increasingly popular, it is important to
understand not only whether they improve answer accuracy but also how they
influence reasoning behaviour, execution time, and mathematical
problem-solving performance. This motivates the need for a systematic
evaluation of different Large Language Models across diverse mathematical
problems using structured prompting strategies.

\section{Problem Statement}

Large Language Models (LLMs) have demonstrated strong performance on
mathematical reasoning benchmarks, with most existing studies evaluating
their capabilities using overall answer accuracy. Although accuracy is an
important performance indicator, it provides only a high-level measure of
model performance and offers limited insight into the reasoning process
that leads to a solution. As a result, existing evaluations often fail to
explain the nature of reasoning failures, the stage at which errors occur,
or how different prompting strategies influence mathematical reasoning
behaviour.

Furthermore, while reasoning-based prompting techniques such as
Chain-of-Thought (CoT) and Tree-of-Thought (ToT) have been widely adopted
to improve mathematical reasoning, many studies compare these approaches
using descriptive accuracy metrics alone. Such comparisons provide limited
evidence regarding whether observed differences in performance are
statistically meaningful. In addition, there is relatively little work
that simultaneously analyses reasoning error types, identifies where
errors occur within the reasoning process, and evaluates the effectiveness
of prompting strategies across multiple Large Language Models using a
controlled experimental design.

To address these limitations, this dissertation presents a systematic
study of mathematical reasoning in Large Language Models using a curated
dataset of mathematical problems covering multiple domains and difficulty
levels. Three Large Language Models are evaluated using both
Chain-of-Thought and Tree-of-Thought prompting strategies. In addition to
measuring answer correctness, the study analyses reasoning error types,
error locations, execution time, and model performance, followed by
appropriate statistical analyses to compare the behaviour of different
models and prompting strategies. By moving beyond aggregate accuracy, this
research aims to provide a more comprehensive understanding of the
reasoning capabilities and limitations of contemporary Large Language
Models.

\section{Research Questions}

This dissertation investigates the mathematical reasoning capabilities of
Large Language Models by addressing the following research questions.

\textbf{RQ1:} What are the most common types of mathematical reasoning
errors made by Large Language Models, and does the distribution of error
types differ across models?

This research question examines whether different Large Language Models
exhibit distinct patterns of mathematical reasoning errors. Since both the
evaluated model and the error type are categorical variables, the study
employs the Chi-square test of independence to determine whether the
observed distribution of error categories differs significantly across
models. Where expected cell frequencies are small, the results are
interpreted with caution because the Chi-square approximation becomes less
reliable. The assumptions and limitations of the Chi-square test,
including the presence of small expected cell counts in this study, are
discussed further in the Methodology chapter.

\textbf{RQ2:} At what stage of the reasoning process do mathematical
reasoning errors first occur, and does this differ across models or
prompting strategies?

This research question investigates whether different models or prompting
strategies tend to fail earlier or later during the reasoning process. The
analysis compares the relative position of the first reasoning error
within generated solutions. As the error-position measurements are not
assumed to follow a normal distribution, the Kruskal--Wallis test is
employed to compare error locations across multiple independent groups.

\textbf{RQ3:} Does Tree-of-Thought prompting significantly improve
mathematical reasoning performance compared with Chain-of-Thought
prompting for each Large Language Model?

This research question evaluates whether Tree-of-Thought (ToT) prompting
produces significantly different answer correctness compared with
Chain-of-Thought (CoT) prompting. Since every mathematical problem is
solved using both prompting strategies for the same model, the
observations are paired rather than independent. Therefore, McNemar's
Exact Test is employed to compare the two prompting strategies while
accounting for the paired experimental design.

Together, these research questions investigate complementary aspects of
mathematical reasoning in Large Language Models. RQ1 analyses the types of
reasoning errors produced by different models, RQ2 examines where
reasoning failures occur within the solution process, and RQ3 evaluates
the effectiveness of two reasoning-based prompting strategies.
Collectively, these questions provide a comprehensive assessment of
mathematical reasoning that extends beyond conventional accuracy-based
evaluations.

\section{Thesis Outline}

The remainder of this dissertation is organised as follows.

\textbf{Chapter 2 -- Literature Review} reviews the existing literature on
Large Language Models (LLMs) and mathematical reasoning. It discusses
Chain-of-Thought (CoT) and Tree-of-Thought (ToT) prompting strategies,
benchmark datasets including Hendrycks MATH and the American Invitational
Mathematics Examination (AIME), and recent approaches to automated
evaluation using LLM-as-a-Judge methodologies.

\textbf{Chapter 3 -- Statistical Methodology} presents the research
methodology adopted in this study. It describes the dataset construction
process, model selection, Tree-of-Thought implementation, answer
evaluation pipeline, error classification framework, and the statistical
methods used to address each of the research questions.

\textbf{Chapter 4 -- Application and Case Study} describes the
implementation of the experimental pipeline used to evaluate the
finalised dataset. It provides a detailed walkthrough of the evaluation
dataset, explains the collection of 786 model responses across different
models and prompting strategies, and describes the data quality assurance
procedures applied to ensure the reliability and consistency of the
experimental results.

\textbf{Chapter 5 -- Results} presents the experimental findings and
statistical analyses. The chapter reports the results for each research
question, compares model performance under Chain-of-Thought and
Tree-of-Thought prompting, and presents analyses of execution time,
answer correctness, reasoning errors, and other evaluation metrics.

\textbf{Chapter 6 -- Discussion and Conclusion} interprets the
experimental findings in relation to the research questions and existing
literature. It discusses the limitations of the study, highlights its key
contributions, provides recommendations for future research, and
concludes the dissertation by summarising the overall findings and their
implications.
